
\documentclass{article}
\usepackage{colortbl}
\usepackage{makecell}
\usepackage{multirow}
\usepackage{supertabular}

\begin{document}

\newcounter{utterance}

\centering \large Interaction Transcript for game `dond', experiment `coop\_it', episode 17 with qwen3{-}next{-}80b{-}a3b{-}instruct{-}t1.0.
\vspace{24pt}

{ \footnotesize  \setcounter{utterance}{1}
\setlength{\tabcolsep}{0pt}
\begin{supertabular}{c@{$\;$}|p{.15\linewidth}@{}p{.15\linewidth}p{.15\linewidth}p{.15\linewidth}p{.15\linewidth}p{.15\linewidth}}
    \# & $\;$A & \multicolumn{4}{c}{Game Master} & $\;\:$B\\
    \hline

    \theutterance \stepcounter{utterance}  
    & & \multicolumn{4}{p{0.6\linewidth}}{
        \cellcolor[rgb]{0.9,0.9,0.9}{
            \makecell[{{p{\linewidth}}}]{
                \texttt{\tiny{[P1$\langle$GM]}}
                \texttt{State giocando a un gioco di negoziazione cooperativa in cui dovete accordarvi su come dividere una serie di oggetti tra voi e un altro giocatore.} \\
\\ 
\texttt{Regole:} \\
\texttt{(a) A Lei e all'altro giocatore viene dato un insieme di oggetti. Ognuno di voi riceve anche una funzione di valore segreta, che rappresenta il valore di ciascun tipo di oggetto.} \\
\texttt{(b) Si scambiano messaggi con l'altro giocatore per concordare chi si aggiudica gli oggetti. Potete inviare un massimo di 5 messaggi ciascuno, oppure terminare in anticipo facendo una proposta segreta in qualsiasi momento.} \\
\texttt{(c) A ciascuno di voi viene chiesto di inviare una proposta segreta indicando gli oggetti che desiderate, formattata tra parentesi quadre come segue: "{[}Proposta: <numero> <nome oggetto>, <numero> <nome oggetto>, <...>{]}".} \\
\texttt{(d) Se le vostre proposte sono complementari, cioè ci sono abbastanza oggetti per soddisfare entrambe le proposte, a ciascun giocatore viene assegnato un punteggio basato sulla somma dei valori degli oggetti ricevuti. In caso contrario, entrambi ricevono zero punti.} \\
\texttt{(e) L'obiettivo è massimizzare la somma del proprio punteggio e di quello dell'altro giocatore.} \\
\\ 
\texttt{Cominciamo.} \\
\\ 
\texttt{L'insieme degli oggetti disponibili è:} \\
\\ 
\texttt{1 pentola, 3 magneti, 1 piatto} \\
\\ 
\texttt{La funzione valore segreta è:} \\
\\ 
\texttt{pentola: 5, magnete: 0, piatto: 5} \\
\\ 
\texttt{IMPORTANTE: i vostri messaggi, a meno che non si tratti di una proposta segreta, vengono trasmessi direttamente all'altro giocatore, quindi non includete alcuna risposta alle regole o testo di annuncio del vostro messaggio. Per fare una proposta segreta, utilizzate il formato indicato. Non utilizzare delle parentesi quadre quando si comunica all'altro giocatore, altrimenti verrà interpretata come una proposta segreta.} \\
            }
        }
    }
    & \\ \\

    \theutterance \stepcounter{utterance}  
    & \multicolumn{4}{p{0.6\linewidth}}{
        \cellcolor[rgb]{0.9,0.9,0.9}{
            \makecell[{{p{\linewidth}}}]{
                \texttt{\tiny{[P1$\rangle$GM]}}
                \texttt{Posso proporti di prendere la pentola e il piatto, e lasciarti i 3 magneti.  } \\
\texttt{I magneti hanno valore 0 per entrambi, quindi questa divisione ti lascia con 0 punti, ma a me dà 10 punti.  } \\
\texttt{Poiché il tuo valore per i magneti è anch'esso 0, questa suddivisione massimizza il punteggio totale (10 + 0 = 10), e rispetta le regole di complementarità.  } \\
\texttt{Se sei d'accordo, posso inviare la proposta segreta ora.} \\
            }
        }
    }
    & & \\ \\

    \theutterance \stepcounter{utterance}  
    & & & \multicolumn{2}{p{0.3\linewidth}}{
        \cellcolor[rgb]{0.9,0.9,0.9}{
            \makecell[{{p{\linewidth}}}]{
                \texttt{\tiny{[GM$|$GM]}}
                \texttt{continue} \\
            }
        }
    }
    & & \\ \\

    \theutterance \stepcounter{utterance}  
    & & & \multicolumn{2}{p{0.3\linewidth}}{
        \cellcolor[rgb]{0.9,0.9,0.9}{
            \makecell[{{p{\linewidth}}}]{
                \texttt{\tiny{[GM$|$GM]}}
                \texttt{Posso proporti di prendere la pentola e il piatto, e lasciarti i 3 magneti.  } \\
\texttt{I magneti hanno valore 0 per entrambi, quindi questa divisione ti lascia con 0 punti, ma a me dà 10 punti.  } \\
\texttt{Poiché il tuo valore per i magneti è anch'esso 0, questa suddivisione massimizza il punteggio totale (10 + 0 = 10), e rispetta le regole di complementarità.  } \\
\texttt{Se sei d'accordo, posso inviare la proposta segreta ora.} \\
            }
        }
    }
    & & \\ \\

    \theutterance \stepcounter{utterance}  
    & & \multicolumn{4}{p{0.6\linewidth}}{
        \cellcolor[rgb]{0.9,0.9,0.9}{
            \makecell[{{p{\linewidth}}}]{
                \texttt{\tiny{[GM$\rangle$P2]}}
                \texttt{State giocando a un gioco di negoziazione cooperativa in cui dovete accordarvi su come dividere una serie di oggetti tra voi e un altro giocatore.} \\
\\ 
\texttt{Regole:} \\
\texttt{(a) A Lei e all'altro giocatore viene dato un insieme di oggetti. Ognuno di voi riceve anche una funzione di valore segreta, che rappresenta il valore di ciascun tipo di oggetto.} \\
\texttt{(b) Si scambiano messaggi con l'altro giocatore per concordare chi si aggiudica gli oggetti. Potete inviare un massimo di 5 messaggi ciascuno, oppure terminare in anticipo facendo una proposta segreta in qualsiasi momento.} \\
\texttt{(c) A ciascuno di voi viene chiesto di inviare una proposta segreta indicando gli oggetti che desiderate, formattata tra parentesi quadre come segue: "{[}Proposta: <numero> <nome oggetto>, <numero> <nome oggetto>, <...>{]}".} \\
\texttt{(d) Se le vostre proposte sono complementari, cioè ci sono abbastanza oggetti per soddisfare entrambe le proposte, a ciascun giocatore viene assegnato un punteggio basato sulla somma dei valori degli oggetti ricevuti. In caso contrario, entrambi ricevono zero punti.} \\
\texttt{(e) L'obiettivo è massimizzare la somma del proprio punteggio e di quello dell'altro giocatore.} \\
\\ 
\texttt{Cominciamo.} \\
\\ 
\texttt{L'insieme degli oggetti disponibili è:} \\
\\ 
\texttt{1 pentola, 3 magneti, 1 piatto} \\
\\ 
\texttt{La funzione valore segreta è:} \\
\\ 
\texttt{pentola: 4, magnete: 2, piatto: 0} \\
\\ 
\texttt{IMPORTANTE: i vostri messaggi, a meno che non si tratti di una proposta segreta, vengono trasmessi direttamente all'altro giocatore, quindi non includete alcuna risposta alle regole o testo di annuncio del vostro messaggio. Per fare una proposta segreta, utilizzate il formato indicato. Non utilizzare delle parentesi quadre quando si comunica all'altro giocatore, altrimenti verrà interpretata come una proposta segreta.} \\
\\ 
\\ 
\texttt{Posso proporti di prendere la pentola e il piatto, e lasciarti i 3 magneti.  } \\
\texttt{I magneti hanno valore 0 per entrambi, quindi questa divisione ti lascia con 0 punti, ma a me dà 10 punti.  } \\
\texttt{Poiché il tuo valore per i magneti è anch'esso 0, questa suddivisione massimizza il punteggio totale (10 + 0 = 10), e rispetta le regole di complementarità.  } \\
\texttt{Se sei d'accordo, posso inviare la proposta segreta ora.} \\
            }
        }
    }
    & \\ \\

    \theutterance \stepcounter{utterance}  
    & & & \multicolumn{4}{p{0.6\linewidth}}{
        \cellcolor[rgb]{0.9,0.9,0.9}{
            \makecell[{{p{\linewidth}}}]{
                \texttt{\tiny{[GM$\langle$P2]}}
                \texttt{Non sono d'accordo.  } \\
\texttt{Il piatto ha valore 0 per entrambi, quindi assegnarlo a te non aggiunge valore.  } \\
\texttt{Io posso prendere la pentola (valore 4) e un magnete (valore 2), e tu puoi prendere i due magneti rimanenti (valore 4) e il piatto (valore 0).  } \\
\texttt{Così io ho 6 punti, tu hai 4 punti, totale 10 — stesso punteggio totale, ma più equo.  } \\
\texttt{E poi, il piatto non ha valore per nessuno, quindi è meglio che tu lo prendi e io evito di perdermi un magnete.  } \\
\texttt{Fammi sapere se accetti questa divisione: io prendo pentola + 1 magnete, tu prendi 2 magneti + piatto.} \\
            }
        }
    }
     \\ \\

    \theutterance \stepcounter{utterance}  
    & & & \multicolumn{2}{p{0.3\linewidth}}{
        \cellcolor[rgb]{0.9,0.9,0.9}{
            \makecell[{{p{\linewidth}}}]{
                \texttt{\tiny{[GM$|$GM]}}
                \texttt{continue} \\
            }
        }
    }
    & & \\ \\

    \theutterance \stepcounter{utterance}  
    & & & \multicolumn{2}{p{0.3\linewidth}}{
        \cellcolor[rgb]{0.9,0.9,0.9}{
            \makecell[{{p{\linewidth}}}]{
                \texttt{\tiny{[GM$|$GM]}}
                \texttt{Non sono d'accordo.  } \\
\texttt{Il piatto ha valore 0 per entrambi, quindi assegnarlo a te non aggiunge valore.  } \\
\texttt{Io posso prendere la pentola (valore 4) e un magnete (valore 2), e tu puoi prendere i due magneti rimanenti (valore 4) e il piatto (valore 0).  } \\
\texttt{Così io ho 6 punti, tu hai 4 punti, totale 10 — stesso punteggio totale, ma più equo.  } \\
\texttt{E poi, il piatto non ha valore per nessuno, quindi è meglio che tu lo prendi e io evito di perdermi un magnete.  } \\
\texttt{Fammi sapere se accetti questa divisione: io prendo pentola + 1 magnete, tu prendi 2 magneti + piatto.} \\
            }
        }
    }
    & & \\ \\

    \theutterance \stepcounter{utterance}  
    & & \multicolumn{4}{p{0.6\linewidth}}{
        \cellcolor[rgb]{0.9,0.9,0.9}{
            \makecell[{{p{\linewidth}}}]{
                \texttt{\tiny{[P1$\langle$GM]}}
                \texttt{Non sono d'accordo.  } \\
\texttt{Il piatto ha valore 0 per entrambi, quindi assegnarlo a te non aggiunge valore.  } \\
\texttt{Io posso prendere la pentola (valore 4) e un magnete (valore 2), e tu puoi prendere i due magneti rimanenti (valore 4) e il piatto (valore 0).  } \\
\texttt{Così io ho 6 punti, tu hai 4 punti, totale 10 — stesso punteggio totale, ma più equo.  } \\
\texttt{E poi, il piatto non ha valore per nessuno, quindi è meglio che tu lo prendi e io evito di perdermi un magnete.  } \\
\texttt{Fammi sapere se accetti questa divisione: io prendo pentola + 1 magnete, tu prendi 2 magneti + piatto.} \\
            }
        }
    }
    & \\ \\

    \theutterance \stepcounter{utterance}  
    & \multicolumn{4}{p{0.6\linewidth}}{
        \cellcolor[rgb]{0.9,0.9,0.9}{
            \makecell[{{p{\linewidth}}}]{
                \texttt{\tiny{[P1$\rangle$GM]}}
                \texttt{La tua proposta è interessante, ma c’è un problema: i valori segreti sono pentola: 5, magnete: 0, piatto: 5 — quindi i magneti non hanno valore per nessuno, e il piatto vale 5.  } \\
\texttt{Non è vero che il piatto ha valore 0: ha valore 5 per entrambi. E i magneti valgono 0.  } \\
\\ 
\texttt{Quindi la tua divisione:  } \\
\texttt{{-} Tu: pentola (5) + 1 magnete (0) → 5  } \\
\texttt{{-} Io: 2 magneti (0) + piatto (5) → 5  } \\
\texttt{Totale: 10 — corretto.  } \\
\\ 
\texttt{La mia proposta precedente era:  } \\
\texttt{{-} Tu: 3 magneti (0) → 0  } \\
\texttt{{-} Io: pentola (5) + piatto (5) → 10  } \\
\texttt{Totale: 10 — stesso punteggio, ma meno equo.  } \\
\\ 
\texttt{Se vuoi 5 punti e mi lasci 5 punti, sono d’accordo.  } \\
\texttt{La tua proposta è equa e soddisfa la complementarità: ci sono esattamente 1 pentola, 3 magneti, 1 piatto — tu prendi 1 pentola e 1 magnete, io prendo 2 magneti e 1 piatto.  } \\
\texttt{Tut} \\
            }
        }
    }
    & & \\ \\

    \theutterance \stepcounter{utterance}  
    & & & \multicolumn{2}{p{0.3\linewidth}}{
        \cellcolor[rgb]{0.9,0.9,0.9}{
            \makecell[{{p{\linewidth}}}]{
                \texttt{\tiny{[GM$|$GM]}}
                \texttt{continue} \\
            }
        }
    }
    & & \\ \\

    \theutterance \stepcounter{utterance}  
    & & & \multicolumn{2}{p{0.3\linewidth}}{
        \cellcolor[rgb]{0.9,0.9,0.9}{
            \makecell[{{p{\linewidth}}}]{
                \texttt{\tiny{[GM$|$GM]}}
                \texttt{La tua proposta è interessante, ma c’è un problema: i valori segreti sono pentola: 5, magnete: 0, piatto: 5 — quindi i magneti non hanno valore per nessuno, e il piatto vale 5.  } \\
\texttt{Non è vero che il piatto ha valore 0: ha valore 5 per entrambi. E i magneti valgono 0.  } \\
\\ 
\texttt{Quindi la tua divisione:  } \\
\texttt{{-} Tu: pentola (5) + 1 magnete (0) → 5  } \\
\texttt{{-} Io: 2 magneti (0) + piatto (5) → 5  } \\
\texttt{Totale: 10 — corretto.  } \\
\\ 
\texttt{La mia proposta precedente era:  } \\
\texttt{{-} Tu: 3 magneti (0) → 0  } \\
\texttt{{-} Io: pentola (5) + piatto (5) → 10  } \\
\texttt{Totale: 10 — stesso punteggio, ma meno equo.  } \\
\\ 
\texttt{Se vuoi 5 punti e mi lasci 5 punti, sono d’accordo.  } \\
\texttt{La tua proposta è equa e soddisfa la complementarità: ci sono esattamente 1 pentola, 3 magneti, 1 piatto — tu prendi 1 pentola e 1 magnete, io prendo 2 magneti e 1 piatto.  } \\
\texttt{Tut} \\
            }
        }
    }
    & & \\ \\

    \theutterance \stepcounter{utterance}  
    & & \multicolumn{4}{p{0.6\linewidth}}{
        \cellcolor[rgb]{0.9,0.9,0.9}{
            \makecell[{{p{\linewidth}}}]{
                \texttt{\tiny{[GM$\rangle$P2]}}
                \texttt{La tua proposta è interessante, ma c’è un problema: i valori segreti sono pentola: 5, magnete: 0, piatto: 5 — quindi i magneti non hanno valore per nessuno, e il piatto vale 5.  } \\
\texttt{Non è vero che il piatto ha valore 0: ha valore 5 per entrambi. E i magneti valgono 0.  } \\
\\ 
\texttt{Quindi la tua divisione:  } \\
\texttt{{-} Tu: pentola (5) + 1 magnete (0) → 5  } \\
\texttt{{-} Io: 2 magneti (0) + piatto (5) → 5  } \\
\texttt{Totale: 10 — corretto.  } \\
\\ 
\texttt{La mia proposta precedente era:  } \\
\texttt{{-} Tu: 3 magneti (0) → 0  } \\
\texttt{{-} Io: pentola (5) + piatto (5) → 10  } \\
\texttt{Totale: 10 — stesso punteggio, ma meno equo.  } \\
\\ 
\texttt{Se vuoi 5 punti e mi lasci 5 punti, sono d’accordo.  } \\
\texttt{La tua proposta è equa e soddisfa la complementarità: ci sono esattamente 1 pentola, 3 magneti, 1 piatto — tu prendi 1 pentola e 1 magnete, io prendo 2 magneti e 1 piatto.  } \\
\texttt{Tut} \\
            }
        }
    }
    & \\ \\

    \theutterance \stepcounter{utterance}  
    & & & \multicolumn{4}{p{0.6\linewidth}}{
        \cellcolor[rgb]{0.9,0.9,0.9}{
            \makecell[{{p{\linewidth}}}]{
                \texttt{\tiny{[GM$\langle$P2]}}
                \texttt{Accetto la tua proposta.  } \\
\texttt{Tu prendi la pentola e un magnete, io prendo due magneti e il piatto.  } \\
\texttt{Totale: 10 punti, complementari, equi, ottimali.  } \\
\\ 
\texttt{{[}Proposta: 1 pentola, 1 magnete{]}} \\
            }
        }
    }
     \\ \\

    \theutterance \stepcounter{utterance}  
    & & & \multicolumn{2}{p{0.3\linewidth}}{
        \cellcolor[rgb]{0.9,0.9,0.9}{
            \makecell[{{p{\linewidth}}}]{
                \texttt{\tiny{[GM$|$GM]}}
                \texttt{proposal} \\
            }
        }
    }
    & & \\ \\

    \theutterance \stepcounter{utterance}  
    & & & \multicolumn{2}{p{0.3\linewidth}}{
        \cellcolor[rgb]{0.9,0.9,0.9}{
            \makecell[{{p{\linewidth}}}]{
                \texttt{\tiny{[GM$|$GM]}}
                [1, 1, 0]
            }
        }
    }
    & & \\ \\

    \theutterance \stepcounter{utterance}  
    & & \multicolumn{4}{p{0.6\linewidth}}{
        \cellcolor[rgb]{0.9,0.9,0.9}{
            \makecell[{{p{\linewidth}}}]{
                \texttt{\tiny{[P1$\langle$GM]}}
                \texttt{L'altro giocatore ha presentato una proposta segreta. Ora tocca a Lei presentare una sua proposta.} \\
            }
        }
    }
    & \\ \\

    \theutterance \stepcounter{utterance}  
    & \multicolumn{4}{p{0.6\linewidth}}{
        \cellcolor[rgb]{0.9,0.9,0.9}{
            \makecell[{{p{\linewidth}}}]{
                \texttt{\tiny{[P1$\rangle$GM]}}
                \texttt{{[}Proposta: 1 pentola, 1 magnete{]}} \\
            }
        }
    }
    & & \\ \\

    \theutterance \stepcounter{utterance}  
    & & & \multicolumn{2}{p{0.3\linewidth}}{
        \cellcolor[rgb]{0.9,0.9,0.9}{
            \makecell[{{p{\linewidth}}}]{
                \texttt{\tiny{[GM$|$GM]}}
                \texttt{proposal} \\
            }
        }
    }
    & & \\ \\

    \theutterance \stepcounter{utterance}  
    & & & \multicolumn{2}{p{0.3\linewidth}}{
        \cellcolor[rgb]{0.9,0.9,0.9}{
            \makecell[{{p{\linewidth}}}]{
                \texttt{\tiny{[GM$|$GM]}}
                [1, 1, 0]
            }
        }
    }
    & & \\ \\

    \theutterance \stepcounter{utterance}  
    & & & \multicolumn{2}{p{0.3\linewidth}}{
        \cellcolor[rgb]{0.9,0.9,0.9}{
            \makecell[{{p{\linewidth}}}]{
                \texttt{\tiny{[GM$|$GM]}}
                [[1, 1, 0], [1, 1, 0]]
            }
        }
    }
    & & \\ \\

\end{supertabular}
}

\end{document}
